\subsection{What the Published Inputs Were}
\label{sec:results:inputs}

Table~\ref{tab:lineage} compares the inputs recovered from the Phase 4A batch file with each candidate series, and Figure~\ref{fig:series} shows the 2008--2026 history (Figure~S3 of the Supplementary Materials shows one window under four conventions).
% src: batch file date 2026-01-04: planning/forensics/2026-10-28_memo_published_inputs.md and scripts/validation/electronics/README.md line 10; windows.csv.gz series_key=jan2026_phase4a (count, min and max window_end); series_comparison.csv pair jan2026_phase4a_vs_new_full scope=all n_days; report.yaml diagnostics.flagged_windows (jan2026_phase4a n_missing_master_days); missing dates: series parquet jan2026_phase4a_SPY against new_full_SPY
The batch file of 2026-01-04 stores the daily values that were sent to o4-mini. They give 1{,}476 windows with ends from 2020-02-13 to 2025-12-31, cut from 1{,}505 daily values. Compared with the complete series, the January calendar lacks 2023-07-12 and 2024-03-11. The 58 windows that span one of the two gaps are flagged by the master calendar, and 87 windows by the NYSE calendar (the 58 and the 29 that span 2020-06-23); all are kept (Section~\ref{sec:methods:windows}, Table~S12 of the Supplementary Materials).

% src: series_comparison.csv pairs jan2026_phase4a_vs_new_full and call_only_vs_jan2026_phase4a, exact_match_pct by year (Table lineage); CALL_FROM and the isolated days: series parquet jan2026_phase4a_SPY, call_only_SPY and new_full_SPY at the 0.01B precision (tolerance 0.005B)
The published inputs follow two different series (Table~\ref{tab:lineage}). In 2020 to 2022 the share of daily values that equal \ident{new_full} at the 0.01B precision of the prompt is 100.0\%, 99.2\% and 98.4\%, and the call-only series equals the inputs on at most 1.2\% of the days. In 2024 and 2025 the relation is reversed: none of the values equals \ident{new_full} (0.0\% in both years) and every value equals the call-only GEX (100.0\% in both years). The year 2023 is the transition: 60.6\% of its values equal \ident{new_full} and 39.4\% equal the call-only series, which together account for every day. From 2023-09-29 on, every value equals the call-only GEX; the day before, 2023-09-28, is the last on which a value differs from it. Before that date, 39 days (the first on 2021-09-17) already equal the call-only value and not the complete one (next paragraph). The published inputs of 2024 and 2025 are therefore the call leg of the options chain without the put leg. They differ in level and sign from the complete series (Table~S7 of the Supplementary Materials gives the means): the mean $|\mathrm{GEX}|$ of the published inputs is \$20.19~billion in 2024 against \$6.48~billion on the complete series, and \$19.21~billion against \$7.13~billion in 2025, with the same daily sign on 62.2\% and 54.0\% of the days.

% src: results/report.yaml sections.published_input_gaps.summary (definition; n_affected_days 39, first_affected_day 2021-09-17, last_affected_day 2023-09-27, affected_days_by_year 2021: 2, 2022: 3, 2023: 34; n_published_days_before_break 940, n_equal_complete_before_break 900, n_equal_neither_before_break 1 (2022-05-31: published -1.33, complete -1.227519 billion); n_days_from_break 565 = n_days_from_break_equal_call_only, n_days_from_break_equal_complete 0; n_recomputed_verdicts_not_equal_to_stored 0 of n_published_windows 1,476); published_input_gaps.csv (39 rows; sign_changes true on 25 of 39); published_input_gaps_by_year.csv scope=all (n_windows 911, n_windows_with_affected_day 332, n_verdicts_changed 46, n_changed_to_detect 45, n_changed_to_not_detect 1)
The 39 earlier call-only days need a definition and their effect. An affected day is a trading day before 2023-09-29 on which the published value equals the stored call-only value (to within \$0.006~billion, the 0.01B rounding of the prompt) and differs from the complete value. Of the 940 published days before that date, 900 equal the complete series, 39 are affected (2 in 2021, 3 in 2022 and 34 in 2023, from 2021-09-17 to 2023-09-27), and one, 2022-05-31, equals neither series: the published value is $-$\$1.33~billion against $-$\$1.23~billion on the complete series, and we leave it as published. The 39 affected days are listed with the three values in \path{results/published_input_gaps.csv}; on 25 of them the complete value has the other sign. We did not establish why these days and no others before 2023-09-29 carry the call leg only. Table~\ref{tab:gaps} shows their effect on the 911 published windows that end before 2023-09-29. With no day replaced, the rule reproduces the stored verdict in all 1{,}476 published windows. The affected days lie in 332 of the 911 windows (36.4\%). Replacing them by the complete values, rounded to the 0.01B of the prompt, changes 46 verdicts of the rule, 45 of them from not detected to detected: the rate for 2022 rises from 32.3\% to 40.6\%, and that for 2023 before the break from 1.6\% to 14.6\%. In every year the rate with the days replaced equals the rule's rate on the complete series for the same end dates. The published inputs of 2021 to 2023 are therefore not a clean reference either, and Table~\ref{tab:replay} should be read with this in mind. The shift of Section~\ref{sec:results:replay} is computed on the complete series and does not depend on these days.

\begin{table}[H]
\centering
\caption{Effect of the call-only days before 2023-09-29 on the rule's verdicts on the published windows that end before that date, SPY (windows by year of the window end; for 2023, ends up to 2023-09-28). Rates are in percent with the windows detected in brackets. An affected day is a trading day on which the published value equals the call-only value within \$0.006~billion and differs from the complete value (Section~\ref{sec:results:inputs}); ``replaced'' sets these days to the complete value, and ``changed'' counts the windows whose verdict differs between the two columns.}
\label{tab:gaps}
\footnotesize
\setlength{\tabcolsep}{4pt}
\begin{tabular}{@{}lrrrrr@{}}
\toprule
Year & Windows & With affected day & Rule, published & Rule, replaced & Changed \\
\midrule
% src: published_input_gaps_by_year.csv scope=year year=2020 (n_windows, n_windows_with_affected_day, n_detect_published, rate_published_pct, n_detect_replaced, rate_replaced_pct, n_verdicts_changed)
2020 & 223 & 0 & 12.1 (27) & 12.1 (27) & 0 \\
% src: published_input_gaps_by_year.csv scope=year year=2021 (same columns)
2021 & 252 & 39 & 4.0 (10) & 3.6 (9) & 1 \\
% src: published_input_gaps_by_year.csv scope=year year=2022 (same columns)
2022 & 251 & 110 & 32.3 (81) & 40.6 (102) & 21 \\
% src: published_input_gaps_by_year.csv scope=year year=2023 (same columns; window_end_last 2023-09-28)
2023 (to 09-28) & 185 & 183 & 1.6 (3) & 14.6 (27) & 24 \\
\midrule
% src: published_input_gaps_by_year.csv scope=all (same columns)
All & 911 & 332 & 13.3 (121) & 18.1 (165) & 46 \\
\bottomrule
\end{tabular}
\end{table}

% src: series_comparison.csv pair old_bs_vs_jan2026_phase4a (exact_match_pct, sign_agree_pct, mean_abs_a_bn) years 2020-2025; missing dates: series parquet old_bs_SPY against new_full_SPY (53 days: 19 spot-filled days of 2009 to 2018, 2023-07-12, 2024-03-11, 26 days 2026-01-02 to 2026-02-09, 2026-05-06, 5 days 2026-10-05 to 2026-10-09)
The pre-rebuild archive (\ident{old_bs}) was computed from the same put-less table for the same dates, so it cannot serve as a clean earlier convention. Its Black--Scholes gamma differs from the stored gamma, and none of its values equals a published value at 0.01B (0.0\% in every year). Its signs agree with the published inputs on 100.0\% of the days in 2024 and 99.6\% in 2025, and its mean $|\mathrm{GEX}|$ is \$16.19~billion and \$14.71~billion, which is the level of a call-only series. The archive also has holes: it lacks 53 days of the complete calendar (the 19 days of 2009 to 2018 and 2026-05-06 that were added on 2026-11-03, 2023-07-12, 2024-03-11, the 26 days from 2026-01-02 to 2026-02-09, and the 5 days from 2026-10-05 to 2026-10-09).

% src: series parquet p3_2024_SPY and p4_2020_SPY (252 values each; sign counts, positive-value range, mean absolute value); windows.csv.gz series_key=p3_2024 and p4_2020 (223 windows each)
The inputs of the Phase 3 and 4 runs of November 2025 survive in the same way: 252 daily values for 2024 and 252 for 2020, 223 windows each, with ISO dates in the prompts. The 2024 values are a gross call-plus-put magnitude with a negative sign imposed. They are negative on 242 days, and the 10 positive values (\$10.61~billion to \$13.76~billion) are placeholder rows of the source database. The 2020 values are a magnitude only and are positive on every day. The mean absolute daily value is \$31.24~billion for 2024 and \$2.85~billion for 2020, so the two baselines differ in convention and in scale (Section~\ref{sec:results:dated}).

\begin{table}[H]
\centering
\caption{Agreement of the recovered published inputs with each candidate series, by year (share of daily values equal at the 0.01B precision of the prompt, in percent). The archive uses another gamma model, so exact equality is not expected for it.}
\label{tab:lineage}
\footnotesize
\begin{tabular}{@{}lrrrr@{}}
\toprule
Year & Days & \ident{new_full} & \ident{call_only} & \ident{old_bs} \\
\midrule
% src: series_comparison.csv pairs jan2026_phase4a_vs_new_full, call_only_vs_jan2026_phase4a, old_bs_vs_jan2026_phase4a, year=2020, n_days and exact_match_pct
2020 & 252 & 100.0 & 0.0 & 0.0 \\
% src: series_comparison.csv pairs jan2026_phase4a_vs_new_full, call_only_vs_jan2026_phase4a, old_bs_vs_jan2026_phase4a, year=2021, n_days and exact_match_pct
2021 & 252 & 99.2 & 0.8 & 0.0 \\
% src: series_comparison.csv pairs jan2026_phase4a_vs_new_full, call_only_vs_jan2026_phase4a, old_bs_vs_jan2026_phase4a, year=2022, n_days and exact_match_pct
2022 & 251 & 98.4 & 1.2 & 0.0 \\
% src: series_comparison.csv pairs jan2026_phase4a_vs_new_full, call_only_vs_jan2026_phase4a, old_bs_vs_jan2026_phase4a, year=2023, n_days and exact_match_pct
2023 & 249 & 60.6 & 39.4 & 0.0 \\
% src: series_comparison.csv pairs jan2026_phase4a_vs_new_full, call_only_vs_jan2026_phase4a, old_bs_vs_jan2026_phase4a, year=2024, n_days and exact_match_pct
2024 & 251 & 0.0 & 100.0 & 0.0 \\
% src: series_comparison.csv pairs jan2026_phase4a_vs_new_full, call_only_vs_jan2026_phase4a, old_bs_vs_jan2026_phase4a, year=2025, n_days and exact_match_pct
2025 & 250 & 0.0 & 100.0 & 0.0 \\
\bottomrule
\end{tabular}
\end{table}



\begin{figure}[H]
\centering
\resultfigure{\figSeries}
% src: series_comparison.csv SPY by year (pairs jan2026_phase4a_vs_new_full, call_only_vs_jan2026_phase4a, new_full_vs_new_ex0dte: mean_abs_a_bn, mean_abs_b_bn, sign_agree_pct)
\caption{SPY, 2008--2026. (a) Yearly mean of $|\mathrm{GEX}|$ under each series, with the \$5 billion magnitude criterion for scale. (b) Daily sign agreement between pairs of series, by year. \textbf{Read this figure as:} until 2022 the published inputs and the complete series are close; from 2023 the published inputs rise to \$20.19~billion in 2024 while the complete series stays between \$6.05~billion and \$7.13~billion in 2023 to 2025. Their daily sign agreement falls from 98.8\% in 2022 to 81.5\%, 62.2\% and 54.0\%. The call-only series matches the published inputs in 2024 and 2025 only. The mean in (a) is taken for each series over the days it shares with one other series of the comparison table, so the day sets differ slightly between series (Table~S7 of the Supplementary Materials); (b) compares each pair on the days that both series hold.}
\label{fig:series}
\end{figure}
